\section{Anexo: extractos de código}
\label{anexo:codigo}

Se incluyen, como \emph{evidencia} de que cada objeto matemático corresponde a código real,
dos extractos del núcleo en C. El código completo está en el repositorio
(\texttt{src/}); el mapeo completa la Tabla de la Sección~\ref{sec:implementacion}.

\subsection{Algoritmo 2 — cálculo de \texorpdfstring{$\tilde k_i$}{k tilde} en \texorpdfstring{$O(M)$}{O(M)}}
Recorre los sitios del cliente del más cercano al más lejano (matriz $S$) acumulando $\bar y$
hasta cubrir $1$ y contando los saltos de radio. Archivo \texttt{src/separation.c}.
\begin{lstlisting}[language=C]
int separation_k_tilde(const SortSites *ss, int i, const double *ybar) {
    int M = ss->M;
    const int  *S = &ss->S[(size_t)i * M];
    const long *D = &ss->Dord[(size_t)i * M];
    int k_tilde = 0, r = 0;
    double val = 1.0 - ybar[S[0]];
    while (val > SEP_EPS && r < M - 1) {
        if (D[r + 1] > D[r]) k_tilde++;   /* salto de radio */
        r++;
        val -= ybar[S[r]];
    }
    return k_tilde;
}
\end{lstlisting}

\subsection{Fase 2 — warm-start: herencia de los cortes de Fase 1}
El maestro MIP precarga el \emph{pool} de cortes de Fase 1 como restricciones antes de ramificar.
Esto es transferencia de información entre fases (warm-start), \emph{no} preprocesamiento de la
matriz $S$. Archivo \texttt{src/phase2.c}.
\begin{lstlisting}[language=C]
/* warm-start: precargar el pool de cortes de Fase 1 como restricciones normales */
if (warm && warm->n > 0) {
    for (int c = 0; c < warm->n; c++) {
        const Cut *cu = &warm->cuts[c];
        wi[0] = M + cu->client; wv[0] = 1.0;            /* theta_i */
        for (int t = 0; t < cu->len; t++) {
            wi[t+1] = cu->ind[t]; wv[t+1] = cu->val[t];
        }
        solver_add_constr(s, cu->len + 1, wi, wv, '>', cu->rhs);
        nwarm++;
    }
}
\end{lstlisting}
